%% kapa: exam=1 round=20 year=2009 period=1 form=A subject=economics subject_ko=경제원론 seq=23 number=63
%% source: Q-Net 감정평가사 기출문제 (cst00309) artlSeq=5024178 file=extracted/1008/5024178/2009년 제20회 감정평가사 1차시험 기출문제/1교시.hwp
%% conversion: kordoc 4.24.51 -> tools/gen1.py (auto); review=proofread-sonnet; flags=none
\begin{question}{63}
케인즈의 이자율이론에 관한 설명으로 옳지 \uline{않은} 것은?
\begin{choices}
  \item 소득수준이 상승하면 화폐수요가 증가한다.
  \item 화폐수요가 증가하면 이자율이 상승한다.
  \item 통화당국이 화폐공급을 증대시키면 이자율이 하락한다.
  \item 통화량이 증가하면 총수요가 감소한다.
  \item 이자율은 화폐의 수요와 공급이 일치하도록 변동한다.
\end{choices}
\end{question}
\answer{4}
